%*************************************
% ภาคผนวก ฉ.
%*************************************
\vspace{-0.5in}

\appendixChapterTitle{ก}{การเตรียม Workspace สำหรับโครงงานแบบหลายรีโพซิทอรี}
\section{ภาพรวมของโครงสร้างแบบหลายรีโพซิทอรี}
\hspace*{1.3em}
โครงงาน Cloud-Based Platform ในปัจจุบันไม่ได้จัดเก็บซอร์สโค้ดแบบรีโพซิทอรีเดียว แต่แยกออกเป็นหลายรีโพซิทอรีตามขอบเขตความรับผิดชอบของบริการเพื่อให้พัฒนา ทดสอบ และติดตั้งใช้งานได้อย่างยืดหยุ่นมากขึ้น รีโพซิทอรีหลักที่ใช้ในโครงงานประกอบด้วย Midori, Momoi, Yuzu และ Yuuka

\begin{mycustomenum2}
    \item \textbf{Midori} เป็นส่วนติดต่อผู้ใช้ พัฒนาด้วย Next.js และ React
    \item \textbf{Momoi} เป็นบริการ API หลัก พัฒนาด้วย Bun, Elysia.js และ Prisma
    \item \textbf{Yuzu} เป็น worker service สำหรับงานจัดการเครื่องเสมือนและงานคิว
    \item \textbf{Yuuka} เป็นรีโพซิทอรีสำหรับไฟล์ deploy, ตัวอย่าง environment, observability และโครงสร้างพื้นฐานที่เกี่ยวข้อง
\end{mycustomenum2}

\section{ข้อกำหนดของระบบ}
\hspace*{1.3em}
ก่อนเริ่มต้นใช้งาน ควรติดตั้งซอฟต์แวร์พื้นฐานให้พร้อมสำหรับการทำงานร่วมกันของทุกรีโพซิทอรีดังนี้

\begin{mycustomenum2}
    \item Bun เวอร์ชัน 1.2 หรือสูงกว่า
    \item Docker และ Docker Compose Plugin
    \item Git สำหรับโคลนซอร์สโค้ดแต่ละรีโพซิทอรี
    \item PostgreSQL และ Redis ในกรณีที่ไม่ใช้งานผ่าน Docker Compose
    \item Proxmox VE สำหรับการทดสอบงานที่เกี่ยวข้องกับการจัดการเครื่องเสมือนจริง
\end{mycustomenum2}

\section{การโคลนรีโพซิทอรีทั้งหมด}
\hspace*{1.3em}
เนื่องจากระบบแยกเป็นหลายรีโพซิทอรี ผู้พัฒนาควรสร้างโฟลเดอร์หลักหนึ่งตำแหน่ง แล้วโคลนแต่ละรีโพซิทอรีลงในโฟลเดอร์เดียวกัน ดังภาพที่ \ref{figF:CloneRepos}

\begin{figure}[htbp]
\centering
\adjustbox{frame, width=0.95\textwidth}{
    \begin{minipage}{0.9\textwidth}
        \texttt{\footnotesize
        mkdir fitm-cloud-platform\\
        cd fitm-cloud-platform\\
        git clone https://github.com/akikungz/midori.git\\
        git clone https://github.com/akikungz/momoi.git\\
        git clone https://github.com/akikungz/yuzu.git\\
        git clone https://github.com/akikungz/yuuka.git
        }
    \end{minipage}
}
\caption{\fontSixTeen{ตัวอย่างคำสั่งสำหรับโคลนรีโพซิทอรีของระบบ}}
\label{figF:CloneRepos}
\end{figure}

\hspace*{1.3em}
เมื่อโคลนครบแล้ว โครงสร้างของ workspace จะประกอบด้วยหลายโฟลเดอร์ระดับเดียวกัน เช่น \texttt{midori}, \texttt{momoi}, \texttt{yuzu} และ \texttt{yuuka} ซึ่งสะท้อนการแบ่งบริการตามหน้าที่ของระบบอย่างชัดเจน

\section{การติดตั้ง dependencies ของแต่ละรีโพซิทอรี}
\hspace*{1.3em}
หลังจากโคลนซอร์สโค้ดแล้ว ต้องติดตั้ง dependencies แยกตามแต่ละรีโพซิทอรี ดังภาพที่ \ref{figF:InstallDeps}

\begin{figure}[htbp]
\centering
\adjustbox{frame, width=0.95\textwidth}{
    \begin{minipage}{0.9\textwidth}
        \texttt{\footnotesize
        cd midori \&\& bun install\\
        cd ../momoi \&\& bun install\\
        cd ../yuzu \&\& bun install
        }
    \end{minipage}
}
\caption{\fontSixTeen{การติดตั้ง dependencies ของแต่ละบริการ}}
\label{figF:InstallDeps}
\end{figure}

\section{การเตรียมไฟล์ Environment}
\hspace*{1.3em}
ค่ากำหนดของระบบในสภาพแวดล้อมพัฒนาถูกแยกเก็บตามรีโพซิทอรี โดยไฟล์ตัวอย่างสำหรับบริการต่าง ๆ ถูกจัดเก็บไว้ในรีโพซิทอรี Yuuka ภายใต้โฟลเดอร์ docs/env ผู้พัฒนาสามารถใช้ไฟล์ตัวอย่างเหล่านี้เพื่อสร้างไฟล์ .env ในรีโพซิทอรีที่เกี่ยวข้องได้

\begin{mycustomenum2}
    \item ใช้ตัวอย่าง midori.env.example สำหรับ Midori
    \item ใช้ตัวอย่าง momoi.env.example สำหรับ Momoi
    \item ใช้ตัวอย่าง yuzu.env.example สำหรับ Yuzu
    \item ตรวจสอบค่า database, redis, storage และ endpoint ต่าง ๆ ให้สอดคล้องกันทุกรีโพซิทอรี
\end{mycustomenum2}

\section{ข้อสังเกตสำหรับการพัฒนาใน workspace เดียวกัน}
\hspace*{1.3em}
แม้ระบบจะเป็นหลายรีโพซิทอรี แต่สามารถเปิดใช้งานร่วมกันใน workspace เดียวของโปรแกรมพัฒนาได้ ซึ่งช่วยให้ค้นหาโค้ด อ้างอิงไฟล์ และติดตามการไหลของข้อมูลระหว่างบริการได้สะดวกขึ้น โดยเฉพาะเมื่อต้องตรวจสอบความสอดคล้องของ schema, API contract และการตั้งค่าของ infrastructure
